\documentclass[11pt,a4paper]{article}
\usepackage[T1]{fontenc}
\usepackage[margin=2.5cm]{geometry}

\pagestyle{empty}

\begin{document}
    
    \begin{center}
        {\Large \textbf{Quasi-Linear Parameter-Varying and Koopman Model Predictive Control for Wind Energy Systems}}\\[0.5cm]
        Antje Dittmer
    \end{center}
    
    \section*{Abstract}
    
    This thesis designs Model Predictive Control algorithms for wind turbines and farms to minimize mechanical loads and power fluctuations. Quasi-linear parameter-varying models enable a qLMPC scheme, applied for the first time as a MIMO controller across the full operating range, which outperforms PI collective pitch control, Individual Pitch Control (IPC), and IPC combined with trailing edge flaps. For wind farm control, Koopman-based linear representations capture wind farm dynamics, enabling adaptive axial induction control to address wind speed variations and combined axial induction and wake redirection control for challenging power setpoints. All designs require minimal measurement availability and are real-time capable.
        
    \section*{Zusammenfassung}
    
Diese Arbeit entwickelt modellprädiktive Regelungsalgorithmen für Windenergieanlagen und -parks, um mechanische Lasten und Leistungsschwankungen zu minimieren. Quasi-lineare parameterveränderliche Modelle ermöglichen ein qLMPC-Verfahren, das zum ersten Mal als Mehrgrößenregler über den gesamten Betriebsbereich eingesetzt wird und eine höhere Regelgüte als ein PI-Regler mit kollektiver Blattverstellung, eine Einzelblattverstellung (IPC) sowie IPC kombiniert mit Hinterkantenklappen erzielt. Für die Windparkregelung wird die Parkdynamik durch Koopman-basierte lineare Modelle repräsentiert. Diese werden sowohl in einer adaptiven Schubbeeinflussungs-Regelung (Axial Induction Control, AIC) eingesetzt, die effizient auf schnelle Windgeschwindigkeitsveränderungen reagiert, als auch in einer kombinierten AIC und Nachlaufumlenkungsregelung (Wake Redirection Control, WRC), die herausfordernde Leistungsvorgaben adressiert. Alle Regelungskonzepte benötigen nur eine geringe Anzahl an Messungen und sind echtzeitfähig.
    
\end{document}